% Local editorial additions to VI/27 (Integral Schemes and Function Fields).
% These examples do not replace canonical CP-VI problem statements or solutions.
\par\medskip
\noindent\textbf{Worked examples: integral schemes, generic points and rational functions.}
An affine scheme $\operatorname{Spec}A$ is integral precisely when $A$ is an
integral domain.  Its unique generic point is $(0)$, and its function field
is $K(X)=\operatorname{Frac}(A)$.  On a non-affine integral scheme the same
field is obtained from any nonempty affine open subset.  The computations
below distinguish an integral total space from the behavior of its fibers,
and a dominant morphism from an isomorphism.

\par\medskip
\noindent\textbf{Worked example 1 (an integral cuspidal curve and its function field).}
Let $k$ be a field and set
\[
 A=k[x,y]/(y^2-x^3),\qquad X=\operatorname{Spec}A.
\]
The assignment $x\mapsto t^2$, $y\mapsto t^3$ defines a homomorphism
$A\to k[t]$.  It is injective: reducing a polynomial modulo
$y^2-x^3$ gives $f(x)+y g(x)$, which maps to
$f(t^2)+t^3g(t^2)$; the first summand has even degrees and the second
has odd degrees, so cancellation is impossible unless both vanish.
Hence $A\cong k[t^2,t^3]\subset k[t]$ and is a domain; therefore $X$
is integral with generic point $(0)$.  In its fraction field,
\[
 t=\frac{t^3}{t^2}=\frac{y}{x},\qquad
 K(X)=\operatorname{Frac}(k[t^2,t^3])=k(t).
\]
The fraction $y/x$ is a rational function on $X$.  It need not be a
global regular function, since $t\notin k[t^2,t^3]$.  This is the basic
difference between $K(X)$ and $\Gamma(X,\mathcal O_X)$.

\par\medskip
\noindent\textbf{Worked example 2 (two components: why a reducible scheme has no single function field).}
Take $B=k[x,y]/(xy)$ and $Y=\operatorname{Spec}B$.
The elements $\bar x,\bar y$ are nonzero but $\bar x\bar y=0$, so $B$
is not a domain.  The ideals $(\bar x)$ and $(\bar y)$ are its two
minimal prime ideals, corresponding to the $y$-axis and $x$-axis.
Each irreducible component has its own function field:
\[
 K\bigl(V(\bar x)\bigr)=k(y),\qquad
 K\bigl(V(\bar y)\bigr)=k(x).
\]
More precisely, the total ring of fractions (localizing at all
non-zero-divisors) is $k(x)\times k(y)$.  Thus it is incorrect to call
$\operatorname{Frac}(B)$ a field: ordinary field-of-fractions notation
requires a domain.  The intersection point is not a third generic point;
its maximal ideal $(\bar x,\bar y)$ lies on both components.

\par\medskip
\noindent\textbf{Worked example 3 (the function field of the projective line).}
Cover $X=\mathbb P^1_k$ by the standard affine charts
\[
 U_0=\operatorname{Spec}k[t],\qquad
 U_\infty=\operatorname{Spec}k[u],\qquad u=t^{-1}
 \text{ on }U_0\cap U_\infty.
\]
Both charts are integral and their overlap is
$\operatorname{Spec}k[t,t^{-1}]$.  Consequently the two fraction fields
identify via $u\mapsto1/t$:
\[
 K(\mathbb P^1_k)=k(t)=k(u).
\]
Now compare regular functions.  A global section is a pair
$f(t)\in k[t]$, $g(u)\in k[u]$ agreeing on the overlap, so
$f(t)=g(t^{-1})$ as Laurent polynomials.  A Laurent polynomial with
only nonnegative powers that also has only nonpositive powers must
be constant.  Therefore
$\Gamma(\mathbb P^1_k,\mathcal O)=k$, even though
$K(\mathbb P^1_k)=k(t)$.
This calculation shows why the function field is not generally the
fraction field of global regular functions on a non-affine scheme.

\par\medskip
\noindent\textbf{Worked example 4 (an integral arithmetic scheme with non-reduced fibers).}
Let $A=\mathbb Z[z]/(z^2-2)$ and $X=\operatorname{Spec}A$.
The polynomial $z^2-2$ is irreducible over $\mathbb Q$, and the monic
presentation injects $A$ into $\mathbb Q(\sqrt2)$.
Thus $A\cong\mathbb Z[\sqrt2]$ is a domain,
$X$ is integral, and $K(X)=\mathbb Q(\sqrt2)$.
For a rational prime $p$, the fiber over $(p)\in\operatorname{Spec}\mathbb Z$
is
\[
 X_{(p)}=\operatorname{Spec}\bigl(\mathbb F_p[z]/(z^2-2)\bigr).
\]
At $p=2$ it is $\operatorname{Spec}(\mathbb F_2[z]/(z^2))$, one
non-reduced point.  At $p=3$, $z^2+1$ has no root in $\mathbb F_3$,
so the fiber is $\operatorname{Spec}\mathbb F_9$, one point with
residue degree two.  At $p=7$,
$z^2-2=(z-3)(z+3)$ with distinct roots, and the fiber is two reduced
$\mathbb F_7$-points.  An integral total space can therefore have
non-reduced or disconnected fibers; integrality does not pass to
arbitrary fibers.

\par\medskip
\noindent\textbf{Worked example 5 (a dominant birational morphism that is not an isomorphism).}
Let $f:\mathbb A^2_k\to\mathbb A^2_k$ correspond to the map
\[
 k[s,t]\longrightarrow k[u,v],\qquad s\longmapsto u,
 \quad t\longmapsto uv.
\]
The map is injective: if $P(u,uv)=0$ as a polynomial in $u,v$,
viewing it as a polynomial in $v$ shows that every coefficient
of $P$ is zero.  Hence $f$ is dominant.  On function fields it induces
\[
 k(s,t)\xrightarrow{\sim}k(u,v),\qquad
 u=s,\quad v=t/s,
\]
so $f$ is birational.  On the open subset $D(s)$ this is an actual
isomorphism, with inverse $(s,t)\mapsto(s,t/s)$.
Nevertheless, if $u=0$, necessarily $(s,t)=(0,0)$; no point with
$s=0$ and $t\ne0$ belongs to the image.  Moreover, the entire
line $u=0$ maps to the origin.  Thus an isomorphism of function fields
proves birationality, not that the underlying morphism is an
isomorphism everywhere.
